\chapter{Answers to Selected Exercises}
\label{app:answers}

Answers are provided for odd-numbered exercises in each chapter. Full
solutions to even-numbered exercises are available in the instructor
supplement. Numerical answers are rounded to three significant figures
unless an exact form is more natural.

\section*{Chapter 0}
\textbf{1.} The only admissible sequence of length $3$ beginning at $a$
is $a \to b \to c$. This matches the adjacency-matrix method: with
vertices ordered $(a,b,c)$, the matrix $M^2$ has a single nonzero entry
in row $a$, at column $c$, confirming exactly one walk of length $2$ from
$a$, ending at $c$.

\section*{Chapter 1}
\textbf{1.} The triangles are similar (SSS: $10/5 = 24/12 = 26/13 = 2$),
with scale factor $k = 2$. \quad
\textbf{3.} The building is approximately $27.5$ meters tall
($h = 34(1.7)/2.1 \approx 27.5$ m). Parallel sunlight is essential
because it guarantees the person-and-shadow triangle and the
building-and-shadow triangle share the same angle of elevation, which is
what makes the two triangles similar in the first place.

\section*{Chapter 2}
\textbf{1.} (a) $5$, (b) $13$, (c) $17$. \quad
\textbf{3.} $\mathbf{v} = \langle 15, 20 \rangle$, $|\mathbf{v}| = 25$ km.
\quad
\textbf{5.} Multiplying by $k$ scales the magnitude by $|k|$. The
direction is unchanged when $k > 0$, and reversed ($180^\circ$ rotation)
when $k < 0$; at $k=0$ the result is the zero vector, which has no
defined direction.

\section*{Chapter 3}
\textbf{1.} Opposite leg $\approx 8.62$ cm, adjacent leg $\approx 11.0$
cm ($14\sin38^\circ$ and $14\cos38^\circ$). \quad
\textbf{3.} The tower is approximately $51.2$ meters tall ($40\tan
52^\circ$).

\section*{Chapter 4}
\textbf{1.} $150^\circ$: reference angle $30^\circ$, Quadrant~II.
$-60^\circ$: reference angle $60^\circ$, Quadrant~IV.
$300^\circ$: reference angle $60^\circ$, Quadrant~IV.
$5\pi/4$ ($=225^\circ$): reference angle $\pi/4$ ($45^\circ$),
Quadrant~III. \quad
\textbf{3.} Three angles coterminal with $40^\circ$ are $400^\circ$,
$-320^\circ$, and $760^\circ$, obtained by adding or subtracting integer
multiples of $360^\circ$; all four angles name the same point of the
unit circle because the coterminal-angle theorem of Chapter~4 guarantees
that angles differing by a multiple of $360^\circ$ produce identical
terminal points.

\section*{Chapter 5}
\textbf{1.} Arc length $s = 15\pi/2 \approx 23.6$ cm; sector area $A =
135\pi/4 \approx 106$ cm$^2$. \quad
\textbf{3.} $x(t) = \cos\left(\dfrac{\pi t}{3} + \dfrac{\pi}{4}\right)$,
$y(t) = \sin\left(\dfrac{\pi t}{3} + \dfrac{\pi}{4}\right)$; at $t=2$ the
angle is $11\pi/12$, giving position approximately $(-0.966, 0.259)$.

\section*{Chapter 6}
\textbf{1.} Amplitude $4$, period $4\pi$, phase shift $2\pi$ to the
right, vertical shift $2$. \quad
\textbf{3.} $\theta = \dfrac{3\pi}{4},\ \dfrac{5\pi}{4},\ \dfrac{11\pi}{4},\
\dfrac{13\pi}{4}$; cosine's graph crosses the value $-\sqrt2/2$ twice in
every period of $2\pi$ (once on the way down, once on the way back up),
and the interval $[0,4\pi)$ spans two full periods, giving four
solutions in all.

\section*{Chapter 7}
\textbf{1.} $\pi/3$. \quad
\textbf{3.} $\pi/4$. \quad
\textbf{5.} $5\pi/6$. \quad
\textbf{7.} $\cos\theta = 15/17$, $\tan\theta = 8/15$ (from the
$8$-$15$-$17$ right triangle). \quad
\textbf{9.} $\theta = \arctan(0.65) \approx 33.0^\circ$.

\section*{Chapter 8}
\textbf{1.} $\sin(105^\circ) = \dfrac{\sqrt6+\sqrt2}{4}$. \quad
\textbf{3.} $\sin(15^\circ) = \dfrac{\sqrt6-\sqrt2}{4}$. \quad
\textbf{5.} $\sin(A-B) = \sin(A+(-B)) = \sin A\cos(-B) + \cos A\sin(-B) =
\sin A\cos B - \cos A\sin B$, using $\cos(-B)=\cos B$ and
$\sin(-B)=-\sin B$ from Chapter~6. \quad
\textbf{7.} Applying the rotation formula, $R(40^\circ)R(50^\circ)$
sends $(1,0)$ to $(\cos90^\circ,\sin90^\circ) = (0,1)$, which is exactly
where $R(90^\circ)$ alone sends $(1,0)$, confirming the two successive
rotations agree with a single $90^\circ$ rotation. \quad
\textbf{9.} $R(90^\circ) = \begin{pmatrix} 0 & -1 \\ 1 & 0
\end{pmatrix}$; applied to $(1,0)$ this gives $(0,1)$.

\section*{Chapter 9}
\textbf{1.} Dividing $\sin^2\theta+\cos^2\theta=1$ by $\cos^2\theta$
gives $\tan^2\theta + 1 = \sec^2\theta$. \quad
\textbf{3.} $\dfrac{\sin\theta}{\cos\theta} = \tan\theta$ directly by the
quotient identity's definition. \quad
\textbf{5.} $\dfrac{\cos^2\theta}{1-\sin\theta} =
\dfrac{(1-\sin\theta)(1+\sin\theta)}{1-\sin\theta} = 1+\sin\theta$,
using $\cos^2\theta = 1-\sin^2\theta = (1-\sin\theta)(1+\sin\theta)$.
\quad
\textbf{7.} $\sin(-\theta)+\cos(-\theta) = -\sin\theta+\cos\theta$, using
the odd symmetry of sine and the even symmetry of cosine. \quad
\textbf{9.} Because $\sin\theta$ and $\cos\theta$ can never both exceed
$1$ in absolute value: the Pythagorean identity forces $\sin^2\theta +
\cos^2\theta = 1$, but $2^2+3^2 = 13 \neq 1$, so no angle $\theta$ can
produce that pair.

\section*{Chapter 10}
\textbf{1.} $\cos(2\theta) = \cos(\theta+\theta) = \cos\theta\cos\theta -
\sin\theta\sin\theta = \cos^2\theta - \sin^2\theta$. \quad
\textbf{3.} Substituting $\cos^2\theta = 1-\sin^2\theta$ into
$\cos(2\theta)=\cos^2\theta-\sin^2\theta$ gives $\cos(2\theta) =
(1-\sin^2\theta)-\sin^2\theta = 1-2\sin^2\theta$. \quad
\textbf{5.} With $\cos\theta=-3/5$ in Quadrant~II, $\sin\theta = 4/5$,
so $\cos(2\theta) = 1-2\sin^2\theta = 1 - 2(16/25) = -7/25$. \quad
\textbf{7.} $\cos(22.5^\circ) = \sqrt{(1+\cos45^\circ)/2} =
\dfrac{\sqrt{2+\sqrt2}}{2}$. \quad
\textbf{9.} $\cos A\cos B = \dfrac{\cos(A+B)+\cos(A-B)}{2}$, obtained by
adding the cosine sum and difference formulas and dividing by $2$.

\section*{Chapter 11}
\textbf{1.} $\theta = \dfrac{\pi}{3}+2\pi k$ or $\theta =
\dfrac{2\pi}{3}+2\pi k$. \quad
\textbf{3.} $\theta = \dfrac{\pi}{4}+\pi k$. \quad
\textbf{5.} $\sin\theta = 0$ gives $\theta = \pi k$; $2\cos\theta+1=0$
gives $\cos\theta = -1/2$, so $\theta = \dfrac{2\pi}{3}+2\pi k$ or
$\theta = \dfrac{4\pi}{3}+2\pi k$. \quad
\textbf{7.} Using $\cos(2\theta)=2\cos^2\theta-1$, the equation becomes
$2\cos^2\theta-\cos\theta-1=0$; substituting $u=\cos\theta$ and
factoring, $(2u+1)(u-1)=0$, so $\cos\theta=1$ (giving $\theta=2\pi k$) or
$\cos\theta=-1/2$ (giving $\theta = 2\pi/3+2\pi k$ or $\theta =
4\pi/3+2\pi k$). \quad
\textbf{9.} Because sine is not one-to-one on all of $\mathbb{R}$:
$\arcsin(a)$ returns only the one solution lying in
$[-\pi/2,\pi/2]$, while infinitely many other angles, differing by the
Quadrant~II reflection and by multiples of $2\pi$, also satisfy
$\sin\theta=a$.
